\eqnrefs{eq:delta-r-definition-dp11} 与\eqnrefs{eq:delta-rho-top-dp11} 之间还可以再向可观测量推进一步。把一圈修正按物理来源组织时，$\Delta r$ 的主导结构可写成
\begin{equation}
\boxed{
\Delta r
=\Delta\alpha
-\frac{c_W^2}{s_W^2}\Delta\rho
+\Delta r_{\rm rem},}
\label{eq:sm-deltar-decomposition-dp95}
\end{equation}
其中 $\Delta\alpha$ 收集从 Thomson 极限到电弱尺度的光子真空极化，$\Delta\rho$ 收集守护对称性破缺的自能差，$\Delta r_{\rm rem}$ 则包含剩余的自能、顶角、盒图以及有限反项组合。这个分解不是把每张图单独赋予物理意义；它只是把完整、规范无关的 $\Delta r$ 按具有清楚极限行为的有限组合重新整理。

把\eqnrefs{eq:delta-r-definition-dp11} 与在壳关系 $s_W^2=1-M_W^2/M_Z^2$ 联立，可写成
\begin{equation}
\boxed{
M_W^2\left(1-\frac{M_W^2}{M_Z^2}\right)
=\frac{\pi\alpha}{\sqrt2G_F^{(E)}}(1+\Delta r).}
\label{eq:sm-mw-deltar-relation-dp95}
\end{equation}
因此固定输入 $(\alpha,G_F^{(E)},M_Z)$ 后，$M_W$ 不再是自由参数，而是对完整电弱圈修正的预测量。在线性阶令 $M_W^2=M_{W,0}^2+\delta M_W^2$，得到
\begin{equation}
\boxed{
\frac{\delta M_W}{M_W}
=-\frac{s_W^2}{2(c_W^2-s_W^2)}\Delta r
+\mathcal O(\Delta r^2).}
\label{eq:sm-mw-linear-response-dp95}
\end{equation}
特别地，只抽出守护对称性破缺项，\eqnrefs{eq:sm-deltar-decomposition-dp95} 给出
\begin{equation}
\boxed{
\left.\frac{\delta M_W}{M_W}\right|_{\Delta\rho}
=\frac{c_W^2}{2(c_W^2-s_W^2)}\Delta\rho.}
\label{eq:sm-mw-rho-response-dp95}
\end{equation}
于是重顶极限不再只是“$\Delta\rho$ 随 $M_t^2$ 增长”的抽象结论，而是直接转化为
\begin{equation}
\left.\frac{\delta M_W}{M_W}\right|_{t}
=\frac{3c_W^2G_F^{(E)}M_t^2}
{16\pi^2\sqrt2\,(c_W^2-s_W^2)}
+\mathcal O\!\left(\frac{M_b^2}{M_t^2}\ln\frac{M_t^2}{M_b^2}\right).
\label{eq:sm-mw-top-response-dp95}
\end{equation}
这条链把 Yukawa 质量劈裂、规范玻色子自能和实验可测的 $W$ 质量放在同一个定量关系里：$m_t$ 通过 $y_t$ 打破守护对称性，自能差形成 $\Delta\rho$，$\Delta\rho$ 再进入 $\Delta r$ 并移动由 $(\alpha,G_F,M_Z)$ 预测的 $M_W$。
